\chapter{Волновое уравнение}

\section{Свойства сферических средних}

$ u \in \mathcal C(\R^n) $

Определим функцию (это среднее значение по сфере вокруг точки $ x $):
$$ I(x, r) = \frac1{|S_r|} \int\limits_{S_r} u(x + y) \di S_y \undereq{y = rz}
\frac1{|S_1|} \int\limits_{S_1} u(x + rz) \di S_z $$
Понятно, что функция непрерывна по $ r $.

\begin{statement}
	$ I(x, 0) = u(x) $
\end{statement}

\begin{proof}
	Совершенно очевидно из второй формулы.
\end{proof}

\begin{statement}
	$ u \in \mathcal C^2(\R^n) $

	$$ \int\limits_{\mathtt B_R} u(x + y) \di y =
	\int_0^R \int\limits_{S_r} u(x + y) \di S_r(y) \di r =
	\int_0^R |S_1| r^{n - 1} I(x, r) \di r $$
	Возьмём оператор Лапласа от левой и правой частей:
	$$ \int_0^R |S_1| r^{n - 1} \Delta_x I(x, r) \di r =
	\int\limits_{\mathtt B_R} \Delta u(x + y) \di y = A $$
	Применим формулу Гаусса"--~Остроградского ($ \Delta u = \dive (\mathcal D u) $):
	\begin{multline*}
		A = \int\limits_{S_R} \bigl( \mathcal D u(x + y), \vv n_y \bigr) \di S_y =
		\int\limits_{S_R} \sum_i \mathcal D_i u(x + y) \cdot \vv n_i(y) \di S_y
		\undereq{y = Rz} \\
		= \int\limits_{S_1} R^{n - 1} \sum_i \mathcal D_i u(x + Rz) \vv n_i(Rz) \di S_z =
		\int\limits_{S_1} R^{n - 1} \sum_i \mathcal D_i (x + Rz) z_i \di S_z = \\
		= \int\limits_{S_1} R^{n - 1} \pder R u(x + Rz) \di S_z = |S_1| R^{n - 1} \pder R  I(x, R)
	\end{multline*}

	Сократим на $ |S_1| $ и продифференцируем по $ R $:
	$$ R^{n - 1} \Delta_x I(x, R) = \pder R \Bigl( R^{n - 1} \pder R I(x, R) \Bigr) $$
	$$ \Delta_x I(x, R) = \frac1{R^{n - 1}} \pder R \Bigl( R^{n - 1} \pder R I(x, R) \Bigr) $$
\end{statement}

\begin{statement}
	$ n = 3 $

	$$ \Delta_x I(x, R) = \frac{\partial^2}{\partial R^2} I(x, R) + \frac2R I(x, R) $$

	Введём функцию $ I_1(x, R) = R I(x, R) $.
	Домножим тождество на $ R $:
	$$ \Delta_x I_1(x, R) = R \frac{\partial^2}{\partial R^2} I(x, R) + 2I(x, R) =
	\frac{\partial^2}{\partial R^2} I_1(x, R) $$
\end{statement}

\section{Решение задачи Коши для волнового уравнения при \texorpdfstring{$ n = 3 $}{n = 3}}

\subsection{Решение задачи \rom1}

В прошлый раз мы выяснили, что достаточно решить задачу \rom1:
$$
\begin{cases}
	u_{tt} - a^2 \Delta u = 0, \quad x \in \R^3, \quad t \ge 0 \\
	u\bigr|_{t = 0} = 0 \\
	u_t\bigr|_{t = 0} = \psi(x)
\end{cases} $$

Возьмём сферическое среднее по $ x $ ($ t $ будет параметром):
$$ I(x, t, R) = \frac1{|S_R|} \int\limits_{S_R} u(x + y, t) \di S_y $$
$$ \frac1{|S_R|} \int\limits_{S_R} u_{tt} (x + Ry, t) \di S_y =
\frac{\partial^2}{\partial t^2} I(x, t, R) $$
$$ \frac1{|S_R|} \int\limits_{S_R} \Delta_x u(x + Ry, t) \di S_y = \Delta_x I(x, t, R) $$
Получили, что
$$ \frac{\partial^2}{\partial t^2} I - a^2 \Delta_x I = 0 $$
Домножим на $ R $:
$$ \frac{\partial^2}{\partial R^2} I_1 - a^2 \Delta_x I_1 = 0 $$
А Лаплас от $ I_1 $ мы уже знаем:
$$ \frac{\partial^2}{\partial R^2} I_1(x, t, R) -
a^2 \frac{\partial^2}{\partial R^2} I_1 (x, t, R) = 0 $$
А это "--- одномерное волновое уравнение ($ x $ теперь играет роль параметра).

При этом, $ t \ge 0 $ и $ \R \ge 0 $.
То есть, мы получили уравнение колебаний полуограниченной струны.
Но нужны ещё два начальных условия и одно граничное условие:
$$ I_1\bigr|_{t = 0} = 0 $$
$$ I_{1t}\bigr|_{t = 0} = R \cdot \frac1{|S_R|} \int\limits_{S_R} \psi(x + y) \di S_y \eqqcolon
\psi_1(x, R) $$
$$ I_1\bigr|_{R = 0} = 0 $$

Предположим, что мы эту задачу решили.
Как теперь вернуться к исходной?
$$ \lim\limits_{R \to 0} \frac{I_1(x, t, R)}R = u(x, t) $$
Воспользуемся формулой Д'Аламбера для $ R \le at $:
$$ I_1(x, t, R) = \frac1{2a} \int_{at - R}^{at + R} \psi_1(x, \rho) \di \rho $$
$$ u(x, t) = \lim\limits_{R \to 0} \frac1{2aR} \int_{at - R}^{at + R} \psi_1(x, \rho) \di \rho =
\frac1a \psi_1(x, at) \bydef t \frac1{|S_1|} \int\limits_{S_1} \psi(x + atz) \di S_z =
\frac{t}{4 \pi} \int\limits_{S_1} \psi(x + atz) \di S $$

\subsection{Решение задачи \rom2}

Напишем решение задачи \rom2:
$$
\begin{cases}
	u_{tt} - a^2 \Delta u = 0 \\
	u\bigr|_{t = 0} = \phi(x) \\
	u_t\bigr|_{t = 0} = \psi(x)
\end{cases} $$
$$ u(x, t) = \pder t \Bigl( \frac t{4\pi} \int\limits_{S_1} \phi(x + atz) \di S_z \Bigr) +
\frac t{4\pi} \int\limits_{S_1} \psi(x + atz) \di S_z $$

Эта формула называется \emph{формулой Кирхгофа}.

Но нужно, чтобы все нужные функции были достаточно дифференцируемы:

\begin{theorem}
	$ \phi \in \mathcal C^3(\R^3), \quad \psi \in \mathcal C^2(\R^3) $

	Тогда $ u \in \mathcal C^2(\R^3 \times \ol \R_+) $, задаваемая уравнением Кирхгофа,
	удовлетворяет однородному волновому уравнению и начальным условиям, то есть, является решением
	задачи Коши.
\end{theorem}

Можно показать, что, если снизить гладкости $ \phi $ и $ \psi $, то функции из $ \mathcal C^2 $ мы
не получим.

\subsection{Физические следствия из формулы Кирхгофа}

Пусть функция $ u $ имеет носитель $ \Omega $ (по $ x $), а мы находимся в точке $ x_0 $
(вне $ \Omega $).
Опуская из некоторой точки $ t $ характеристический конус, получаем некоторый шар в $ X $.
Обозначим за $ t_* $ ту точку, в которой этот шар начинает касаться $ \Omega $, $ t^* $ "--- ту,
в которой перестаёт касаться.

Из формулы понятно, что и до $ t_* $, и после $ t^* $ интегралы равны нулю.

\comment{Тут есть картиночка, можно будет нарисовать, если совсем скучно станет.}

\section{Решение задачи Коши для волнового уравнения при \texorpdfstring{$ n = 2 $}{n = 2}}

\subsection{Решение задачи \rom1}

$$
\begin{cases}
	u_{tt} - a^2 \Delta_2 u = 0, \quad x \in \R^2, \quad t \ge 0 \\
	u\bigr|_{t = 0} = 0 \\
	u_t\bigr|_{t = 0} = \psi(x)
\end{cases} $$

Будем решать другую задачу:
$$
\begin{cases}
	u_{tt} - a^2 \Delta_3 u = 0 \\
	u\bigr|_{t = 0} = 0 \\
	u_t\bigr|_{t = 0} = \psi(x_1, x_2)
\end{cases} $$
\begin{multline*}
	u(x, t) = \frac t{4\pi} \int\limits_{S_1} (x_1 + atz_1, x_2 + atz_2) \di S_z
	\undereq{\int\limits_{S} = 2\int\limits_{\mathtt B}}
	\frac t{2\pi} \int\limits_{z_1^2 + z_2^2 \le 1} \frac{\psi(x_1 + atz_1, x_2 + atz_2)}
	{\sqrt{1 - z_1^2 - z_2^2}} \di z_1 \di z_2 = \\
	= \int\limits_{|y| \le at} \frac{\psi(x + y)}{\sqrt{(at)^2 - |y|^2}} \di y
\end{multline*}
Эту формулу иногда называют \emph{формулой Пуассона"--~Кирхгофа}.

\subsection{Решение задачи \rom2}

$$
\begin{cases}
	u_{tt} - a^2 \Delta_2 u = 0 \\
	u\bigr|_{t = 0} \phi(x) \\
	u_t\bigr|_{t = 0} = \psi(x)
\end{cases} $$
$$ u(x, t) = \pder t \Bigl( \frac1{\pi a} \int\limits_{|y| \le at}
\frac{\phi(x + y)}{\sqrt{(at)^2 - |y|^2}} \di y \Bigr) +
\frac1{2 \pi a} \int\limits_{|y| \le at} \frac{\psi(x + y)}{\sqrt{(at)^2 - |y|^2}} \di y $$
Условия те же, что и трёхмерном случае: $ \phi \in \mathcal C^3(\R^2) $, $ \psi \in
\mathcal C^2(\R^2) $.

\subsection{Физические следствия из формулы Пуассона\texorpdfstring{"--~}{--}Кирхгофа}

Картинка та же: носитель $ \Omega $, характеристический конус и точка $ t_* $.
Но в формуле Пуассона"--~Кирхгофа интегрирование ведётся по кругу, поэтому мы не можем утверждать,
что колебания прекратятся при прохождении через некоторую точку $ t^* $.

Но в какой-то момент расширяющийся шар поглотит $ \Omega $, и все интегралы станут интегралами по
фиксированной области.
А подынтегральное выражение убывает.
Значит, колебания затухают, но не прекращаются.
Это явление называют \emph{диффузией волн}.
Можно показать, что во всех чётных размерностях имеет место диффузия волн, а во всех нечётных "---
задний фронт.

\chapter{Уравнение теплопроводности}

Пусть $ \Omega \sub \R^n $.
Рассмотрим цилиндр $ Q = \Omega \times (0, \tau] $.

$$
\begin{cases}
	u_t - a^2 \Delta u = f(x, t) \text{ в } Q \\
	u\bigr|_{S_{\text{бок}}} = u_0(x, t) \\
	u\bigr|_{t = 0} = \phi(x)
\end{cases} $$
$$ u \in \mathcal C_{x,t}^{2,1} (Q) \cap \mathcal C(\ol Q) $$

\section{Принцип максимума для уравнения теплопроводности}

\begin{statement}
	$ u \in \mathcal C_{x,t}^{2,1}(Q) \cap \mathcal C(\ol Q) $
	$$ u_t - a^2 \Delta u = 0 \text{ в } Q $$

	$$ \implies \max\limits_{\ol Q} u = \max\limits_{\ol Q \setminus Q} u, \quad
	\min\limits_{\ol Q} u = \min\limits_{\ol Q \setminus Q} u $$
\end{statement}

\begin{proof}
	\textbf{Пусть} $ m \coloneq \max\limits_{\ol Q \setminus Q} u <
	M \coloneq \max\limits_{\ol Q} u $.
	$$ \exists (x^0, t^0) \in Q : \quad M = u(x^0, t^0) $$

	Рассмотрим функцию
	$$ u_\eps(x, t) = u(x, t) + \eps (t^0 - t) $$
	$$ u_\eps(x^0, t^0) = M $$
	На $ \ol Q \setminus Q \quad u_\eps \le m + \eps \tau $.
	Выберем $ \eps $ так, чтобы $ m + \eps \tau < M $.

	Функция $ u_\eps $ достигает максимума не на стакане:
	$$ \max e_\eps = e_\eps(x^1, t^1), \quad (x^1, t^1) \in Q $$
	Воспользуемся необходимым условием локального максимума:
	$$ (u_\eps)_t (x^1, t^1) \ge 0 $$
	(для $ t^1 < \tau $ будет равенство)
	$$ \mathcal D_i \mathcal D_i u_\eps (x^1, t^1) \le 0 $$
	$$ \Bigl( (u_\eps)_t - a^2 \Delta u_\eps \Bigr)(x^1, t^1) \ge 0 $$
	Но $ (u_\eps)_t - a^2 \Delta u_\eps \equiv -\eps $ "--- \contra.
\end{proof}

\subsection{Следствия из принципа максимума}

\begin{theorem}[единственности, для начальной краевой задачи]
	Вот та задача не может иметь более, чем одно решение.
\end{theorem}

\begin{proof}
	\textbf{Пусть} имеется два решения, $ v $ "--- их разность.
	$$
	\begin{cases}
		v_t - a^2 \Delta v = 0 \\
		v\bigr|_{S_{\text{бок}}} = 0 \\
		v\bigr|_{t = 0}
	\end{cases} $$
	\textbf{Но} по принципу максимума, максимум находится на стакане.
\end{proof}

\begin{theorem}
	$$
	\begin{cases}
		u_t - a^2 \Delta u = 0 \text{ в } Q = \Omega \times \R_+ \\
		u\bigr|_{S_{\text{бок}}} = 0 \\
		u\bigr|_{t = 0} = \phi(x)
	\end{cases} $$

	$$ \implies u(x, t) \uniarr[t \to \infty]{} 0 $$
\end{theorem}

\begin{proof}
	$$ w(x, t) \coloneq A \cos (cx_1) e^{-bt} $$
	$$ w_t - aw^2 = w(-b + a^2c^2) = 0 \implies b = a^2 c^2 $$
	$$ \Omega \sub \Set{x \in \R^n | {} |x_1| < d} $$
	Хотим, чтобы в этом слое $ \cos $ был положительным.
	Возьмём $ c = \frac{\pi}{3d} $.
	$$ w\bigr|_{S_{\text{бок}}} \ge \frac12 A e^{-bt} \ge |u| = 0 $$
	$$ w\bigr|_{t = 0} \ge \frac12 A $$
	$$ A \coloneq 2 \max |\phi| $$
	Тогда $ w\bigr|_{t = 0} \ge |u| = |\phi| $.

	Теперь $ w $ мажорирует $ |u| $ на $ \ol Q \setminus Q $.
	$$ \implies w \pm u \ge 0 \text{ на } \ol Q \setminus Q $$
	По принципу максимума, $ w \pm u \ge 0 $ в $ \ol Q $.
	$$ \implies w \ge |u| \text{ на } \ol Q $$
	$$ |u(x, t)| \le A e^{-bt} $$
\end{proof}
